\section{Introduction}
\label{sec:introduction}
Inpatient diabetes care involves changing glucose levels, treatment decisions, and measurements taken at intervals that vary across patients and clinical settings. A glucose value of 180 mg/dL, for example, may reflect persistent hyperglycemia or a rapid fall from a much higher level after insulin. Interpreting that value in context matters when estimating a patient's risk of subsequent complications.

Measurements can be interpreted in different ways. Knowledge-based temporal abstraction uses clinical rules to identify states, trends, and events in a patient's history \citep{shahar1997framework}. The same glucose readings, for example, can indicate both a high glucose state and a falling glucose trend. Statistical methods can also group measurements into intervals, using patterns in the data rather than clinical definitions. In the MIMIC-IV intensive care study, INTERVenE-Enc learned from both kinds of intervals and distinguished patients who developed complications more accurately than models of individual observations, including STraTS \citep{oded2026intervene,johnson2023mimiciv}.

Across six Israeli general hospitals in the Kineret network, a clinical study used a knowledge base developed for inpatient diabetes care to describe glucose trajectories and assess adherence to national treatmentrecommendations \citep{oded2027kineret-jbi,hochberg2023inpatient}. Greater adherence to the recommendation to give insulin after hyperglycemia was associated with less severe hyperglycemia, with only a slight increase in the risk of hypoglycemia. The same records allow us to ask a different question. Given what is already known about a patient's glucose levels and treatment, does an explicit assessment of adherence help predict a subsequent complication?

We used the first four days of each eligible Kineret admission to predict complications through day 14, their time of onset, and total length of stay. We examined INTERVenE-Enc in general inpatient care in a second health system, comparing its performance with models of individual measurements and testing what clinical intervals added over statistical intervals in the same encoder. We also assessed whether measures of adherence to inpatient diabetes-care guidelines added predictive information and examined how the encoder's named clinical patterns contributed to risk predictions.
